\documentclass[12pt]{article}
\usepackage[T1]{fontenc}
\usepackage[utf8]{inputenc}
\usepackage{lmodern}
\usepackage{textcomp}
\usepackage{enumitem}
\usepackage{geometry}
\usepackage{graphicx}
\usepackage{amsmath, amssymb}
\geometry{margin=1in}
\begin{document}
\begin{center}
Astronomy - Graduate Level Assessment 4 \
Total Marks: 40 \
Answer all questions.
\end{center}

\section*{Multiple Choice (10 marks)}
\begin{enumerate}[label=\arabic*.]
\item Kirkwood gaps are observed in the main asteroid belt including at the position(s) where:
\begin{enumerate}[label=(\alph*)]
    \item asteroids would orbit with a period half that of Jupiter
    \item asteroids would orbit with a period twice that of Jupiter
    \item asteroids would orbit with a period twice that of Mars
    \item A and B
\end{enumerate}
\item What are the conditions necessary for a terrestrial planet to have a strong magnetic field?
\begin{enumerate}[label=(\alph*)]
    \item fast rotation only
    \item a rocky mantle only
    \item a molten metallic core only
    \item both a molten metallic core and reasonably fast rotation
\end{enumerate}
\item We were first able to accurately measure the diameter of Pluto from:
\begin{enumerate}[label=(\alph*)]
    \item a New Horizons flyby in the 1990 s
    \item Hubble Space Telescope images that resolved Pluto's disk
    \item brightness measurements made during mutual eclipses of Pluto and Charon
    \item radar observations made by the Arecibo telescope
\end{enumerate}
\item Which is not a similarity between Saturn and Jupiter's atmospheres?
\begin{enumerate}[label=(\alph*)]
    \item a composition dominated by hydrogen and helium
    \item the presence of belts zones and storms
    \item an equatorial wind speed of more than 900 miles per hour
    \item significant "shear" between bands of circulation at different latitudes
\end{enumerate}
\item The terrestrial planet cores contain mostly metal because
\begin{enumerate}[label=(\alph*)]
    \item the entire planets are made mostly of metal.
    \item metals condensed first in the solar nebula and the rocks then accreted around them.
    \item metals sank to the center during a time when the interiors were molten throughout.
    \item radioactivity created metals in the core from the decay of uranium.
\end{enumerate}
\end{enumerate}

\section*{True / False (10 marks)}
\textit{Answer True or False}
\begin{enumerate}[label=\arabic*.]
\setcounter{enumi}{5}
\item True or False: Sunspots are black regions on the Sun that experience a significant increase in number every 9 years, referred to as the solar cycle.
\item True or False: The lithosphere consists of the softer rocky material found in the mantle beneath the crust.
\item True or False: Observing a star through cool hydrogen gas reveals only emission lines characteristic of the star's composition.
\item True or False: Volcanic heating cannot account for Mercury's lack of a permanent atmosphere due to its insufficient volcanic activity.
\item True or False: Helium is the second most abundant element in the solar system, following hydrogen in terms of prevalence.
\end{enumerate}

\section*{Long Form Response (20 marks)}
\textit{Answer each question with a clear and complete explanation.}
\begin{enumerate}[label=\arabic*.]
\setcounter{enumi}{10}
\item Discuss the reasons why Saturn's size is comparable to Jupiter's despite having a significantly lower mass, focusing on the effects of mass on density and planetary structure.
\item Discuss the arrangement of electromagnetic radiation categories from shortest to longest wavelength, analyzing their implications in astronomy and how they affect observational techniques.
\end{enumerate}

\vfill
\noindent\textit{End of Paper}
\end{document}